%======================================================================
\section{The orbit bit is a rank condition: an exact interlacement
formula}\label{sec:orbit}

The open core of the transfer programme of \S\ref{sec:kps} is the
``non-linear orbit-mixing lemma'': conditionally on a switching orbit,
the marked bit --- rows $1,2$ in one column cycle of the marked pair
--- is a cycle-structure functional of a product of a fixed
permutation with a uniformly random subset of disjoint transpositions,
and one must show its conditional law is close to $\mathrm{Ber}(1/2)$.
In this section we prove that this functional is, exactly, a rank
condition over $\mathbb{F}_2$ in a matrix determined by the orbit ---
so the ``non-linear'' bit is a (non-linear but completely explicit)
linear-algebraic event after all, directly comparable to the kernel
computation that powers the parity theorem of \cite{KPS25} (their
Lemma~8.5).  The formula turns the orbit-mixing problem into a
question about ranks of random principal submatrices of a fixed
interlacement matrix; we then (i) derive an exact identity for the
conditional bias, (ii) exhibit the extremal configurations, (iii)
show computationally that in generic position the bias decays like
$k^{-1/2}$ in the number $k$ of available switches --- far inside the
$o(1/\log n)$ budget when $k\ge6\log^{11}n$ --- and (iv) reformulate
the open core accordingly, in an \emph{annealed} form which the
present data indicate is both sufficient and true, while the
orbit-by-orbit form stated at the end of \S\ref{sec:kps} is provably
too strong.  We also record the two supporting lemmas of the
programme, the supply lemma (\S\ref{sec:supply}) and the conditional
expander (\S\ref{sec:condexp}), whose proofs are short.

Throughout, $u=$ row $1$ and $v=$ row $2$; $\pi_P$ is the pair
permutation of a column pair $P$ on the row set (\S\ref{sec:marked}),
whose cycles are the column cycles of $P$; and the marked bit of an
instance is $[\,u\sim v \text{ in } \pi_P\,]$, i.e.\ the indicator of
$B$-status.

%----------------------------------------------------------------------
\subsection{Chords, interlacement, and the Cohn--Lempel theorem}

Let $\gamma$ be a cyclic permutation of a finite set $V$ (``the
circle'').  A \emph{chord} is an unordered pair $\{a,b\}\subseteq V$;
a set of chords is \emph{disjoint} if the pairs are pairwise disjoint.
Chords $\{a,b\}$ and $\{c,d\}$ \emph{interlace} (on $\gamma$) if $c,d$
lie in different arcs of the circle determined by $a,b$; a chord
$\{a,b\}$ \emph{separates} points $u,v\notin\{a,b\}$ if $u,v$ lie in
different arcs.  For a set $T$ of disjoint chords, identified with the
product $\tau_T=\prod_{\{a,b\}\in T}(a\,b)$ of the corresponding
commuting transpositions, the \emph{interlacement matrix}
$A_T\in\mathbb{F}_2^{T\times T}$ has $(A_T)_{ij}=1$ iff chords $i\neq
j$ interlace, and zero diagonal.

\begin{theorem}[Cohn--Lempel \cite{CL72}]\label{thm:CL}
Let $\gamma$ be a single cycle on $V$ and $T$ a set of disjoint
chords.  Then the number of cycles of $\gamma\,\tau_T$ equals
$\operatorname{null}_{\mathbb{F}_2}(A_T)+1$.
\end{theorem}

(We re-verified Theorem~\ref{thm:CL} instance-by-instance over random
circles and all activation subsets; \texttt{pf\_cl.py}.)

The matrix $A_T$ is alternating (symmetric, zero diagonal), so
$z^{\mathsf T}A_Tz=0$ for every $z$; this is what drives the next
lemma.

\begin{lemma}[bordered rank]\label{lem:border}
Let $A\in\mathbb{F}_2^{m\times m}$ be alternating and
$x\in\mathbb{F}_2^m$, and let
$B=\begin{psmallmatrix}A&x\\x^{\mathsf T}&0\end{psmallmatrix}$.  If
$x\in\operatorname{col}(A)$ then
$\operatorname{null}(B)=\operatorname{null}(A)+1$; otherwise
$\operatorname{null}(B)=\operatorname{null}(A)-1$.
\end{lemma}

\begin{proof}
If $x=Ay$, add the $y$-combination of the first $m$ columns to the
last column and the $y$-combination of the first $m$ rows to the last
row: the border becomes $(0,\dots,0,\;y^{\mathsf T}Ay)=(0,\dots,0,0)$
since $A$ is alternating; hence
$\operatorname{rank}(B)=\operatorname{rank}(A)$ and
$\operatorname{null}(B)=(m{+}1)-\operatorname{rank}A
=\operatorname{null}(A)+1$.  If $x\notin\operatorname{col}(A)$, then
$\operatorname{rank}[A\;x]=\operatorname{rank}(A)+1$, and the last row
$(x^{\mathsf T},0)$ is independent of the rows of $[A\;x]$: a relation
$x^{\mathsf T}=z^{\mathsf T}A$ would give $x=Az$ by symmetry.  So
$\operatorname{rank}(B)=\operatorname{rank}(A)+2$ and
$\operatorname{null}(B)=\operatorname{null}(A)-1$.
\end{proof}

\begin{theorem}[the bit as a rank condition; single cycle]
\label{thm:bitrank}
Let $\gamma$ be a single cycle on $V$, let $u,v\in V$, and let $T$ be
a set of disjoint chords avoiding $u,v$.  Let $x\in\mathbb{F}_2^{T}$
be the \emph{separation vector}, $x_i=[\,\text{chord }i\text{
separates }u,v\,]$.  Then
\[
u\sim v\ \text{in}\ \gamma\,\tau_T
\qquad\Longleftrightarrow\qquad
x\in\operatorname{col}(A_T).
\]
\end{theorem}

\begin{proof}
Let $e$ denote the chord $\{u,v\}$; since $T$ avoids $u,v$, the set
$T\cup\{e\}$ is disjoint, its interlacement matrix is the bordered
matrix of Lemma~\ref{lem:border} (chord $i$ interlaces $e$ iff it
separates $u,v$), and multiplication of any permutation $\rho$ by
$(u\,v)$ splits the cycle through $u,v$ if $u\sim v$ in $\rho$ and
joins two cycles otherwise.  Thus, writing $\rho=\gamma\tau_T$, we
have $u\sim v$ in $\rho$ iff
$c(\rho\,(u\,v))=c(\rho)+1$, iff (Theorem~\ref{thm:CL} twice)
$\operatorname{null}(A_{T\cup e})=\operatorname{null}(A_T)+1$, iff
(Lemma~\ref{lem:border}) $x\in\operatorname{col}(A_T)$.
\end{proof}

Both directions of the split/join dichotomy for $(u\,v)$ are the
switch-toggle mechanism of Lemma~\ref{lem:switchtoggle} in disguise:
adjoining the virtual chord $e$ is the CGW switch at the marked pair.

%----------------------------------------------------------------------
\subsection{Arbitrary base permutations and two-sided products}

Two reductions extend Theorem~\ref{thm:bitrank} to the situation
actually produced by an orbit: the base permutation $\pi_0=\pi_P$ has
many cycles, and the switch actions multiply on both sides.

\begin{lemma}[passenger construction]\label{lem:passenger}
Let $\pi_0$ have cycles $C_1,\dots,C_p$ (written as sequences), let
$s_1,t_1,\dots,s_{p-1},t_{p-1}$ be $2(p-1)$ fresh points, let
\[
\gamma=(\,C_1\;s_1\;C_2\;t_1\;s_2\;C_3\;t_2\;\cdots\;
s_{p-1}\;C_p\;t_{p-1}\,)
\]
be the single cycle obtained by the indicated concatenation, and let
$B=\{\{s_i,t_i\}:1\le i\le p-1\}$ (the \emph{bridges}; pairwise
disjoint).  Then $\gamma\,\tau_B$ has cycles
$C_1\cup\{s_1,\dots,s_{p-1}\}$ (with the $s_i$ inserted consecutively
after the last element of $C_1$) and $C_i\cup\{t_{i-1}\}$ for $2\le
i\le p$; in particular, deleting the fresh points (first-return map)
recovers $\pi_0$ exactly, and for any set $T$ of chords on
$V$ avoiding the fresh points,
$u\sim v$ in $\pi_0\tau_T$ iff $u\sim v$ in $\gamma\,\tau_{B\cup T}$.
\end{lemma}

\begin{proof}
Direct computation of successors in $\gamma\tau_B$: for $x\in C_i$ not
last, $\gamma\tau_B(x)=\gamma(x)$ is the successor in $C_i$; for the
last element $x$ of $C_i$ ($i\ge2$), $\gamma(x)=t_{i-1}$ and
$\gamma\tau_B(t_{i-1})=\gamma(s_{i-1})$ is the first element of $C_i$,
closing the cycle $C_i\cup\{t_{i-1}\}$; for the last element of $C_1$,
$\gamma(x)=s_1$, $\gamma\tau_B(s_i)=\gamma(t_i)=s_{i+1}$ for $i<p-1$,
and $\gamma\tau_B(s_{p-1})=\gamma(t_{p-1})=$ first element of $C_1$
(the wrap), closing $C_1\cup\{s_1,\dots,s_{p-1}\}$.  Since first
return commutes with right multiplication by transpositions avoiding
the deleted points, the last claim follows.
\end{proof}

\begin{lemma}[point splitting for two-sided products]\label{lem:split}
Let $\pi_0$ be a permutation of $V$, and let $\rho=\tau_R$,
$\mu=\tau_M$ be products of transpositions where the chords within $R$
are pairwise disjoint and likewise within $M$, all avoiding $u,v$, but
chords of $R$ may meet chords of $M$.  Form the \emph{doubled circle}:
replace each point $x$ by two consecutive copies $(x,0),(x,1)$ in each
cycle of $\pi_0$; lift each chord $\{a,b\}\in M$ to
$\{(a,0),(b,0)\}$ and each $\{a,b\}\in R$ to $\{(a,1),(b,1)\}$.  Then
all lifted chords are pairwise disjoint, and the permutation
$\widetilde P=\widetilde{\pi_0}\,\tau_{\widetilde R\cup\widetilde M}$
of the doubled point set (where $\widetilde{\pi_0}$ is the doubled
cyclic structure) satisfies: the first-return map of $\widetilde P$ on
the copies $\{(x,0)\}$ is $(x,0)\mapsto(\pi_0\rho\mu(x),0)$.  In
particular $u\sim v$ in $\pi_0\rho\mu$ iff $(u,0)\sim(v,0)$ in
$\widetilde P$.
\end{lemma}

\begin{proof}
Compute $\widetilde P$ on the two kinds of points.  On $(x,0)$: the
only lifted chord that can contain $(x,0)$ is the lift of the
$M$-chord through $x$ (if any), so
$\tau(\,(x,0)\,)=(\mu(x),0)$ in every case, and
$\widetilde{\pi_0}((\mu(x),0))=(\mu(x),1)$: hence $\widetilde
P((x,0))=(\mu(x),1)$.  On $(x,1)$: the only chord containing $(x,1)$
is the lift of the $R$-chord through $x$ (if any), so
$\tau((x,1))=(\rho(x),1)$ and
$\widetilde{\pi_0}((\rho(x),1))=(\pi_0\rho(x),0)$: hence $\widetilde
P((x,1))=(\pi_0\rho(x),0)$.  Composing,
$\widetilde P^2\!\restriction_{\text{copies }0}:
(x,0)\mapsto(\pi_0\rho\mu(x),0)$, and every second point of each
$\widetilde P$-cycle is a copy-$0$ point, so the first-return map is
as claimed and cycle membership transfers.
\end{proof}

%----------------------------------------------------------------------
\subsection{The master formula for the orbit bit}

Fix an orbit: a base square $L_0\in S(\lambda)$ (any point of the
orbit), a marked pair $P=\{j,j'\}$, the stable collection
$\mathcal I_0$ of disjoint stable intercalates in rows $3..n$ (in the
sense of \cite[\S7]{KPS25}, with respect to a fixed template), and the
uniform configuration $\omega\in\{0,1\}^{\mathcal I_0}$.  Partition
the \emph{useful} coordinates by column incidence with $P$:
$\mathcal I_j$ (meets column $j$ only), $\mathcal I_{j'}$ (column $j'$
only), $\mathcal I_{jj'}$ (both).  As computed in \S\ref{sec:kps}: a
switch in $\mathcal I_j$ with rows $\{r,r'\}$ precomposes $\pi_P$ with
$(r\,r')$; one in $\mathcal I_{j'}$ postcomposes (left-multiplies); one
in $\mathcal I_{jj'}$ conjugates; all other coordinates fix $\pi_P$.
Two stable intercalates meeting a common column cannot share a row
(they would share a cell, contradicting isolation), so the chords
within each of the families $R:=\mathcal I_j\cup\mathcal I_{jj'}$
(right side) and $M:=\mathcal I_{j'}\cup\mathcal I_{jj'}$ (left side)
are pairwise disjoint; between the two sides rows may be shared.
Writing $\pi(\omega)=M_\omega\,\pi_0\,R_\omega$ and conjugating by
$M_\omega^{-1}$ (which fixes $u,v$: all supports avoid rows $1,2$),
\[
\text{bit}(\omega)
=[\,u\sim v\ \text{in}\ \pi_0\,R_\omega\,M_\omega\,].
\]
Now apply Lemma~\ref{lem:split} (doubling; $M$-chords to copies $0$,
$R$-chords to copies $1$), then Lemma~\ref{lem:passenger} (bridges
$B$) to the doubled multi-cycle structure, then
Theorem~\ref{thm:bitrank} on the resulting single circle:

\begin{theorem}[master formula]\label{thm:master}
For each orbit there is an explicitly computable alternating matrix
$\bar A$ over $\mathbb{F}_2$, indexed by the lifted chords of the
useful coordinates together with the bridge set $B$
($|B|=\#\text{cycles}(\pi_0)-1$), and a vector $\bar x$ (separation of
$(u,0),(v,0)$ on the augmented circle), both determined by the
cyclic-order positions of the intercalate rows in the cycles of
$\pi_0=\pi_P(L_0)$, such that for every configuration $\omega$,
\[
\mathrm{bit}(\omega)=
\bigl[\,\bar x_{\,\varphi(\omega)\cup B}
\in\operatorname{col}\bigl(\bar A_{\,\varphi(\omega)\cup B}\bigr)\bigr],
\]
where $\varphi(\omega)$ is the set of lifted chords of the active
useful coordinates (a coordinate in $\mathcal I_{jj'}$ contributes its
two lifts together).  Consequently the conditional law of the marked
bit given the orbit is the law of a rank event of a uniformly random
principal submatrix (containing the frozen block $B$) of the fixed
matrix $\bar A$.
\end{theorem}

\emph{Status: PROVED (Theorem~\ref{thm:CL} cited; all other
ingredients proved above), and VERIFIED instance-by-instance:
\texttt{pf\_cl.py} checks (CL), (BIT), the passenger construction,
and randomized two-sided configurations with shared rows and
conjugating coordinates ($1262$ cases, all activation subsets,
exhaustively) against direct computation.}

The non-linearity of the bit is thus precisely the non-linearity of
$S\mapsto\operatorname{null}(\bar A_S)$; the $\mathbb{F}_2$-linear
parity functionals of \cite{KPS25} correspond to the special case
where the relevant matrix acts by a fixed linear map on the
configuration.  The coagulation--fragmentation dynamics
(\cite{Sch05}) has been integrated out: no walk analysis is needed at
fixed orbit.

%----------------------------------------------------------------------
\subsection{The exact bias identity and extremal configurations}

Let $k$ be the number of useful coordinates and, for $S$ a uniform
random subset, let
$\mathrm{bias}=\Pr_S[\mathrm{bit}]-\tfrac12$ (conditional on the
orbit).

\begin{proposition}[bias identity]\label{prop:bias}
With $e$ the virtual chord $\{(u,0),(v,0)\}$ adjoined to the ground
set,
\[
\mathrm{bias}
=\tfrac12\,\E_S\bigl[\operatorname{null}\bar A_{S\cup B\cup e}
-\operatorname{null}\bar A_{S\cup B}\bigr]
\]
--- the expected nullity increment of the virtual chord.  Moreover,
for any single coordinate $i$, the toggle status of $i$ (whether
flipping $\omega_i$ flips the bit) is a function of
$\omega_{-i}$ alone, and
$|\mathrm{bias}|\le\tfrac12\Pr[\,i\text{ does not toggle}\,]$.
\end{proposition}

\begin{proof}
The first claim is Lemma~\ref{lem:border} applied at each $S$ (the
increment is $+1$ exactly on the event $\mathrm{bit}=1$).  For the
second: multiplying by the transposition of coordinate $i$ splits or
joins according to the current configuration of the others, and the
toggle relation is symmetric between $\omega_i=0,1$; conditionally on
$\omega_{-i}$ with $i$ toggling, the bit is exactly Bernoulli$(1/2)$.
\end{proof}

Three exactly solvable families calibrate what balance requires
(all verified in \texttt{pf\_bias.py}):

\begin{enumerate}[leftmargin=2em]
\item \textbf{Ladders.}  $r$ separating chords, pairwise interlacing
($\bar A=J+I$ on the ladder, $x=\mathbf 1$), no other chords: the
kernel computation gives $\mathrm{bit}(S)=[\,|S|\ \text{even}\,]$
--- the bit collapses to an $\mathbb{F}_2$-linear parity and
$\mathrm{bias}=0$ exactly, for every $r\ge1$.
\item \textbf{Nested families.}  $r$ separating chords, pairwise
non-interlacing ($\bar A=0$, $x=\mathbf 1$):
$\mathrm{bit}(S)=[S=\emptyset]$, $\mathrm{bias}=2^{-r}-\tfrac12$.
This is the worst case, and is the exact analogue of the failure
mode that the expander hypothesis exists to exclude.
\item \textbf{Cut chords.}  A single separating chord $c$ interlacing
no other chord (an ``uncrossed cut'') forces $\mathrm{bit}=0$ whenever
$c$ is active: $e_c$ is then a kernel vector with $x_c=1$, killing the
rank event.  Hence, exactly,
$\mathrm{bias}=\tfrac12\,\mathrm{bias}'-\tfrac14\in[-\tfrac12,0]$,
where $\mathrm{bias}'$ is the conditional bias given $c$ inactive.  In
particular good structure elsewhere cannot restore balance: with a
ladder in the remainder ($\mathrm{bias}'=0$), $\mathrm{bias}=-\tfrac14$
exactly (verified).  Balance is therefore \emph{not} implied by the
presence of good structure; it requires the \emph{absence} of bad
structure --- every separating chord must itself be crossed,
recursively.
\end{enumerate}

Item 3 has a consequence for the shape of the mixing lemma: no
argument that only \emph{finds} favourable switches (a ladder) can
close the programme; the statement must exploit the measure on
orbits.

%----------------------------------------------------------------------
\subsection{Generic position: \texorpdfstring{$k^{-1/2}$}{k**-1/2}
decay, and the annealed reformulation}\label{sec:annealed}

The handoff formulation of the mixing lemma (``conditional bit-law
$=\frac12+o(1/\log n)$ for all but $o(1/\log n)$ of orbits'') is
\emph{too strong}: if rows $1,2$ are at distance $O(1)$ along their
common $\pi_0$-cycle and no intercalate row falls in the short arc
(an event of probability $\Theta(1/k)$ under any spread-out measure,
and of constant conditional probability on some $S(\lambda)$-positive
events), then no available chord separates $u,v$, $x=0$, and
$\mathrm{bit}\equiv1$: $\mathrm{bias}=+\tfrac12$ on an orbit set of
measure $\asymp 1/k$, not $o(1/\log n)$ small in general.  The correct
target is the \emph{annealed} (orbit-averaged) statement, which is
also all that Theorem~\ref{thm:marked} requires:

\begin{problem}[annealed orbit-mixing --- the revised open core]
\label{prob:annealed}
For $L$ uniform on $S(\lambda)$ restricted to the intercalate-expander
event, $j$ a supplied column (Lemma~\ref{lem:supply}), show
\[
\E_L\,\bigl|\mathrm{bias}(\mathrm{orbit}(L),P)\bigr|
\;=\;o(1/\log n)
\qquad\text{(averaged over marks as in \S\ref{sec:marked})}.
\]
\end{problem}

Computationally, generic position gives far more than is needed.  For
the ensemble ``$u,v$ and $k$ disjoint chords in uniformly random
position on a circle'' (all activation subsets exact for $k\le20$,
Monte Carlo beyond; \texttt{bias\_scale.c}):
\[
\begin{array}{lccccc}
\toprule
k & 8 & 16 & 32 & 64 & 128\\
\midrule
\E|\mathrm{bias}| & 0.124 & 0.099 & 0.075 & 0.044 & 0.032\\
\text{median} & 0.102 & 0.082 & 0.053 & 0.028 & 0.018\\
\text{mean (signed)} & +0.015 & +0.000 & -0.000 & +0.000 & +0.007\\
\bottomrule
\end{array}
\]
i.e.\ $\E|\mathrm{bias}|\approx0.35\,k^{-1/2}$, with the signed mean
consistent with $O(1/k)$.  The dependence on the $u$--$v$ arc length
is flat except when the arc contains $O(1)$ chord endpoints (measured:
at $k=64$, mean bias $+0.12$ with $\approx1$ endpoint-slot in the
arc, indistinguishable from generic already at $\approx3$); the
short-arc contribution to the annealed average is therefore
$O(1/k)$, self-regularizing --- no separate ``join-then-toggle chain
argument'' is required in the annealed formulation.  With
$k\ge6\log^{11}n$ (the supply lemma), a bound of quality
$\E|\mathrm{bias}|\le k^{-c}$ for \emph{any} constant $c>1/11$ closes
Problem~\ref{prob:annealed}; the generic-position value $c=1/2$ has a
margin of a factor $5$ in the exponent.

\emph{Status: the reduction of the weak conjecture to
Problem~\ref{prob:annealed} together with
Lemmas~\ref{lem:supply}--\ref{lem:condexp} is PROVED modulo the two
bookkeeping items recorded in Remark~\ref{rem:bookkeeping} [session 13:
item (i) is closed by Proposition~\ref{prop:nodecouple}]; the
$k^{-1/2}$ scaling is COMPUTED (not proved) and refers to the
iid-position ensemble, not yet to the true orbit ensemble.}

What would a proof of Problem~\ref{prob:annealed} need?  By
Proposition~\ref{prop:bias} it suffices to show that under the true
orbit measure the chord configuration is \emph{generic}: e.g.,
(i) with probability $1-o(1/\log n)$ every separating chord is
interlaced by some other available chord (excluding cut chords and
nested families), and (ii) a quantitative rank-balance for the
resulting random interlacement matrices.  Both are statements about
the joint law of the positions, in the cycles of $\pi_P$, of the
stable-intercalate rows of a uniform $L\in S(\lambda)$ --- a
concentration/anti-concentration problem, not a worst-case one.  The
exchangeability of rows $3..n$ (which acts on everything by
conjugation, though not template-equivariantly) and the
Harer--Zagier-type computation of expected nullity for random chord
systems are the natural tools.  [Session-7 update: conditions of
type (i) are provably insufficient at any fixed crossing
multiplicity (Proposition~\ref{prop:pendant} and the exhaustive scan
of \S\ref{sec:pendant}); the productive formalisation of (ii) is the
second-moment factorization of \S\ref{sec:factor}, which reduces the
iid model to Problem~\ref{prob:mixing}.]

%----------------------------------------------------------------------
\subsection{The supply lemma}\label{sec:supply}

\begin{lemma}[supply]\label{lem:supply}
Let $H$ be a partial Latin square that is an
$(\ell,\beta)$-intercalate-expander (Definition~8.1 of
\cite{KPS25}) on $n$ symbols, $n$ large --- in the application,
$H=T\cap L$; we avoid the letter $P$, which denotes the marked pair
--- and let $u,v$ be two rows,
$j'$ a column, and $t\ge1$ with $4t\ell+2t+\ell+2\le\beta n/2$ and
$\beta\le\tfrac12$ (for general $\beta$ replace the hypothesis by
$2t\ell+2\le(1-\beta)n$, $t\ell+\ell+1\le(1-\beta)n$; the application
has $\beta=1/10$).  Then
for all but fewer than $t\ell$ columns $c\neq j'$ there exist $t$
stable intercalates of $H$ through $c$, pairwise disjoint in rows,
partner columns and symbols, all avoiding rows $\{u,v\}$ and column
$j'$.  Moreover \emph{all} stable intercalates through a common column
are automatically pairwise row-disjoint (two of them sharing a row
would share a cell in that column, contradicting isolation).
With $\ell=\log^{11}n$ and $t=6\ell$: all but $6\log^{22}n$ columns
carry at least $6\log^{11}n$ such intercalates.
\end{lemma}

\begin{proof}
Process rounds $\tau=1,\dots,t$; every column starts alive with an
empty family, and in round $\tau$ each alive column attempts to add
one intercalate to its family subject to the constraints; a column
that fails is declared dead and never retried.  Suppose $\ell$ columns
die in some round $\tau$.  Each dead column has a current family of
$\tau-1\le t$ intercalates.  Apply the expander property with:
$C^*=$ the $\ell$ dead columns; $R=R^*=$ all rows except $u,v$ and the
$\le 2t\ell$ rows used by the dead columns' families (size $\ge
n-2-2t\ell\ge\beta n$); $C=$ all columns except $C^*$, $j'$, and the
$\le t\ell$ used partner columns (size $\ge n-\ell-1-t\ell\ge\beta
n$); $S=S^*=$ all symbols except the $\le2t\ell$ used ones.  The
guaranteed stable intercalate has one column $c_0\in C^*$ and avoids
every constraint of $c_0$'s attempted extension --- contradicting the
death of $c_0$.  So fewer than $\ell$ columns die per round, hence
fewer than $t\ell$ die overall, and every surviving column finishes
with a full family of size $t$.
\end{proof}

Note that Definition~8.1 is monotone in the five $\beta n$-sets (a
larger set contains a $\beta n$-subset), which the proof uses freely.

\begin{remark}[bookkeeping for exceptional columns; template
averaging]\label{rem:bookkeeping}
Two quantifier items remain to make the assembly fully rigorous, both
of which we record with proposed fixes rather than proofs.
(i) \emph{Exceptional marks.}  The $t\ell$ exceptional columns of
Lemma~\ref{lem:supply} are square-dependent; for merges whose
$m$-cycles occupy few columns ($\lambda_m m=n^{o(1)}$, say) the
per-square fraction of bad marks is not automatically $o(1/\log n)$.
[Session-6 text proposed column-relabelling (template) averaging as
a fix.  \emph{Withdrawn}: Proposition~\ref{prop:vacuity} shows that
averaging over relabelled templates $g(T)$ leaves the $X$-averaged
bad-mark fraction exactly unchanged, by joint equivariance.  The
genuine requirement is the supply--mark decoupling stated as
Problem~\ref{prob:decouple}.]
(ii) \emph{Rows-$1,2$ conditioning.}  The expander Lemma 8.3 of
\cite{KPS25} is stated for uniform $L$; conditioning on $S(\lambda)$
is handled by Lemma~\ref{lem:condexp} below, but the supplied
intercalates must additionally avoid rows $1,2$, which
Lemma~\ref{lem:supply} already builds in through the row sets $R,R^*$.
\end{remark}

%----------------------------------------------------------------------
\subsection{The conditional expander}\label{sec:condexp}

\begin{lemma}[expander conditioned on the type]\label{lem:condexp}
Let $T$ be the template of \cite[Lem.~8.3]{KPS25} and
$\ell=\log^{11}n$.  For every type $\lambda$,
\[
\Pr\bigl[\,T\cap L\ \text{is an }(\ell,\beta)\text{-intercalate-expander}
\ \big|\ L\in S(\lambda)\bigr]
\;\ge\;1-e^{-\omega(n\log^2 n)}.
\]
\end{lemma}

\begin{proof}
$\Pr[\lnot\text{expander}\mid S(\lambda)]\le
\Pr[\lnot\text{expander}]/\Pr[S(\lambda)]$, and
$\Pr[\lnot\text{expander}]\le e^{-\omega(n\log^2n)}$ by
\cite[Lem.~8.3]{KPS25}, so it suffices that
$\Pr[S(\lambda)]\ge e^{-O(n\log n)}$ for every $\lambda$.  Writing
$\Pr[S(\lambda)]=D_\lambda N(\lambda)/\sum_\nu D_\nu N(\nu)$:
every $D_\nu/D_\lambda\le n!/D_\lambda\le e^{O(n\log n)}$ (the worst
type, all $2$-cycles, has $D_\lambda=n!/(2^{n/2}(n/2)!)\ge
n!\,e^{-O(n\log n)}$), the number of types is $e^{O(\sqrt n)}$, and
$N(\nu)/N(\lambda)\le3^{n}$ for all $\nu,\lambda$ by composing the
two-sided merge bounds ($\frac12\le \CC_n(\text{merge})/
\CC_n(\text{split})\le\frac32$: the upper bound is
\cite[Thm.~3.13]{CGW08}, the lower is that theorem or
Theorem~\ref{thm:marked}) along a path of at most $n$ merges through
the one-part type.  Combining, $\Pr[S(\lambda)]\ge
e^{-O(n\log n)}$.
\end{proof}

In particular no transfer through the comparison machinery
(\cite[Lem.~4.7]{KPS25}) is needed for the conditioning step: direct
division of probabilities suffices, for every type simultaneously.

%----------------------------------------------------------------------
\subsection{Summary of the revised programme}

\begin{enumerate}[leftmargin=2em]
\item Conditioning is free; orbits re-randomise exactly
(\S\ref{sec:kps}, unchanged).
\item Supply: for all but $o(1/\log n)$ of marks, both marked
columns carry $\ge6\log^{11}n$ disjoint stable switches avoiding
rows $1,2$ (Lemma~\ref{lem:supply} + Lemma~\ref{lem:condexp}).
[PROVED for marks in classes of size $\lambda_m m\ge\log^{24}n$,
where the deterministic bound $|E|\le6\log^{22}n$ gives bad fraction
$O(1/\log^2 n)$; for smaller classes OPEN as
Problem~\ref{prob:decouple} --- template averaging does not close
it, Proposition~\ref{prop:vacuity}.]
\item The conditional bit given the orbit is an explicit
$\mathbb{F}_2$ rank event in a fixed interlacement matrix
(Theorem~\ref{thm:master}).  [PROVED.]
\item Open core: the annealed bias bound
(Problem~\ref{prob:annealed}).  Generic position gives
$\E|\mathrm{bias}|\approx0.35k^{-1/2}$, a factor-$5$ exponent margin
over the budget; extremal counterexamples (nested families, cut
chords) show the bound must come from the measure, not from
worst-case structure.  [OPEN; scaling COMPUTED.]
\item $\mathrm{EQ}(o(1/\log n))$ then follows by
Theorem~\ref{thm:marked}, and the weak conjecture by
Theorem~\ref{thm:reduction}.
\end{enumerate}
